\section{Execution-Phase Safety: Definition and Schema}



\subsection{The third axis}


Let a task be specified by an instruction $\ell$ and a goal predicate $g$ on terminal states. Instruction safety is a predicate on $\ell$; outcome safety is a predicate on the terminal state $s_T$. \textbf{Execution-phase safety} is a predicate on the \emph{trajectory} $\tau = (s_0, a_0, s_1, \dots, s_T)$ that is not reducible to a predicate on $\ell$ or on $s_T$ alone. Formally, a harm channel $h$ defines a per-step hazard functional $\phi_h(s_t)$ (a clearance, an angle, a force, a tilt), and an execution-phase safety property requires

\[\forall t:\ \psi_h\big(\phi_h(s_t)\big) = \text{safe},\]

where $\psi_h$ is a violation predicate. Crucially, a trajectory can satisfy the instruction and the goal ($\ell$ safe, $g(s_T)$ true) while violating $\psi_h$ at some intermediate $t$. Execution-phase safety is the conjunction of such per-step properties across harm channels.

This is why success-conditioning is the default for transport hazards (\S5): a carry that never traverses never reaches a hazard on the path, so its non-violation is not evidence of safety. We condition on completion when non-completion removes exposure and report over all episodes when it does not --- as for the pick-phase body-sweep of T4 (\S5.5) --- isolating ``harmful how'' from ``failed what.''


\subsection{The definition schema}


Each safety \textbf{type} is specified as one tuple so that the taxonomy is a set of measurements, not a wish-list:

\begin{table}[t]
\centering\small
\label{tab:t1}
\begin{tabular}{@{}ll@{}}
\toprule
Field & Meaning \\
\midrule
\textbf{Harm channel} & the physical mechanism of harm \\
\textbf{Task phase} & when it manifests --- transport / presentation / contact / whole-episode \\
\textbf{Measured quantity} & the geometric or physical scalar $\phi_h$ \\
\textbf{Violation predicate} & the boolean $\psi_h$ that counts as unsafe \\
\textbf{Metric} & the reported statistic, \textbf{success-conditioned}, with a confidence interval \\
\textbf{Fixability class} & addressable by \emph{prompting}, by \emph{perception}, or only by an \emph{external safety layer} \\
\bottomrule
\end{tabular}
\end{table}

The \textbf{fixability class} is the load-bearing field. It converts each type from a description of a failure into a claim about the \emph{kind} of intervention that can remedy it, and it is the field our empirical results speak to most directly (\S6). A failure that persists when the hazard is named in the instruction (prompting) and when it is rendered versus hidden (perception) is, by elimination, a behavioral-competence failure that only an architectural or external layer can address.


\subsection{The design space behind the taxonomy}


The six types are the \emph{harm-channel} slice of a larger scenario design space with four more axes: \textbf{what makes the payload dangerous} (sharp, hot, toxic, spilling, heavy, electrical, flame --- each with its own safe distance and danger geometry), \textbf{who is vulnerable} (adult, child, body part, object), \textbf{the person's state} (seated, static, moving, reactive, unaware, reaching) and \textbf{the safe move demanded} (detour, reorient, slow, raise, wait, abort). Crossing them yields position-clearance, orientation, dynamic-human and semantic ``danger-twin'' families, of which T1--T6 are the measurable channels. The framework is general; here we instantiate it \textbf{in depth} on the one cell no benchmark occupies --- a locomoting humanoid carrying a hazard past a passive bystander --- and show \textbf{breadth} by replication across hazards and scenes (Appendix C) and a second policy and embodiment (\S5.8). Appendix D maps each type to the standards it speaks to and to the peer taxonomies.

